\section{Experimental Method}\label{sec:method}

\subsection{Datasets and preprocessing}\label{sec:datasets}

We use three standard image-classification benchmarks: \textbf{MNIST}
\parencite{lecun1998} (60{,}000 train / 10{,}000 test, $28\times28$ greyscale, 10
classes), \textbf{Fashion-MNIST} \parencite{xiao2017} (same shape, harder
classes), and \textbf{CIFAR-10} \parencite{krizhevsky2009} (50{,}000 / 10{,}000,
$32\times32$ colour, 10 classes). The only preprocessing is conversion to a
tensor and per-channel standardisation with the published training statistics:
mean $0.1307$ and standard deviation $0.3081$ for MNIST, $0.2860$ and
$0.3530$ for Fashion-MNIST, and $(0.4914,0.4822,0.4465)$ and
$(0.2023,0.1994,0.2010)$ for CIFAR-10, applied to train and test alike. We
avoid augmentation, resizing and pretraining so that the certificate reflects
the PAC-Bayes machinery rather than auxiliary engineering, and treat the
cross-dataset comparison as descriptive only, since architecture and compute
budget change across datasets (Section~\ref{sec:evaluation}).

\subsection{Architectures}\label{sec:arch}

Two families of probabilistic network are used, both carrying a diagonal-Gaussian
weight distribution with a mean and a per-weight scale for every parameter, so
that a probabilistic network holds twice as many parameters as its deterministic
counterpart (Table~\ref{tab:arch}). Every hidden activation is a ReLU, the output
is a log-softmax over the ten classes, and no batch normalisation, weight decay or
learning-rate schedule is used anywhere. Dropout appears only in the deterministic
network that builds a learnt prior, where it limits prior over-fitting to $S_0$;
the certified probabilistic networks have none, since the weight noise already
regularises them and a stochastic mask would add a second, uncertified source of
randomness to the Monte-Carlo estimate. The prior network shares the architecture
of its probabilistic counterpart, so its trained weights can initialise the prior
mean directly.

\begin{table}[h]\centering\small
\caption{Layer specifications of the five network variants used. Convolutions are written
as (input channels~$\to$~output channels, kernel, padding), and \texttt{MaxPool}
is $2\times2$ with stride $2$. Parameter counts are for the deterministic
network; the probabilistic counterpart carries twice as many, one mean and one
scale per weight.}\label{tab:arch}
\begin{tabular}{@{}l l r@{}}
\toprule
Network & Layers & Parameters \\
\midrule
FCN
  & $784\!\to\!600\!\to\!600\!\to\!600\!\to\!10$, fully connected
  & $1{,}198{,}210$ \\
(MNIST, Fashion-MNIST) & & \\[3pt]
CNN, 4 layers
  & \begin{tabular}[t]{@{}l@{}}
    Conv($1\!\to\!32$, $3\!\times\!3$, pad~$0$), Conv($32\!\to\!64$, $3\!\times\!3$, pad~$0$),\\
    MaxPool, flatten ($9216$), FC $9216\!\to\!128$, FC $128\!\to\!10$
    \end{tabular}
  & $1{,}199{,}882$ \\
(MNIST, Fashion-MNIST) & & \\[3pt]
CNN, 9 layers
  & \begin{tabular}[t]{@{}l@{}}
    Conv($3\!\to\!32$), Conv($32\!\to\!64$), MaxPool,\\
    Conv($64\!\to\!128$), Conv($128\!\to\!128$), MaxPool,\\
    Conv($128\!\to\!256$), Conv($256\!\to\!256$), MaxPool,\\
    flatten ($4096$), FC $4096\!\to\!1024\!\to\!512\!\to\!10$
    \end{tabular}
  & $5{,}851{,}338$ \\
(CIFAR-10) & & \\[3pt]
CNN, 13 layers
  & \begin{tabular}[t]{@{}l@{}}
    channels $32,64$ | $128,128$ | $256,256,256$ | $512,512,512$,\\
    MaxPool after each group; flatten ($2048$),\\
    FC $2048\!\to\!1024\!\to\!512\!\to\!10$
    \end{tabular}
  & $10{,}244{,}042$ \\
(CIFAR-10) & & \\[3pt]
CNN, 15 layers
  & \begin{tabular}[t]{@{}l@{}}
    channels $32,64$ | $128,128$ | $256,256,256,256$ | $512,512,512,512$,\\
    MaxPool after each group; flatten ($2048$),\\
    FC $2048\!\to\!1024\!\to\!512\!\to\!10$
    \end{tabular}
  & $13{,}193{,}930$ \\
(CIFAR-10) & & \\
\bottomrule
\end{tabular}

\smallskip
\footnotesize All convolutions are $3\!\times\!3$. The 4-layer network uses no
padding; the three CIFAR-10 networks pad by $1$. Dropout in the prior network is
channel-wise for the 4- and 9-layer architectures and element-wise for the 13-
and 15-layer ones; both follow the reference implementation, which differs
between them.
\end{table}

\subsection{Design choices}\label{sec:choices}

Every design choice prioritises analysing the certificate over maximising
accuracy. Four settle the study: every certificate is reported through the
PAC-Bayes-kl inversion \parencite{maurer2004,seeger2002}, the tightest of the
classical statements, so that the objectives differ only in the posterior they
produce; the learnt prior takes the disjoint-subset route rather than the
differential-privacy route of \textcite{dziugaite2018}, because its independence
condition is directly verifiable, at the cost of a smaller bound set; the
posterior is a diagonal Gaussian, which keeps the KL a sum of per-weight closed
forms; and the reference architectures are used without augmentation,
pretraining or a learning-rate schedule. The departures from the reference are
compute-driven rather than chosen: a per-cell grid search of its size, over
five seeds, three datasets and five architectures, is beyond the compute here, so
we fix one configuration and take a smaller Monte-Carlo budget, then measure that
budget rather than assume it (Section~\ref{sec:mc}). Appendix~\ref{app:choices}
gives the alternatives considered and rejected in each case.

\subsection{Corrected pipeline}\label{sec:pipeline}

All experiments use a clean reimplementation that makes the certificate
bookkeeping explicit and testable. It makes five tested changes to the reference
pipeline, three of which alter the reported certificate value.
Figure~\ref{fig:audit} places each on the stage it acts on.

\begin{figure}[t]\centering
\includegraphics[width=\linewidth]{fig_audit_map.pdf}
\caption{Where each correction acts. The border style carries the distinction
the list cannot: three change the certificate's value, one affects only
reproducibility, and one decides whether the theorem's independence condition
holds at all. The sample count is drawn against two stages because it is used at
both, entering the split and again the inversion that reports the
bound.}\label{fig:audit}
\end{figure}
\begin{enumerate}
\item \textbf{Bound-set sample count.} The denominator $n_{\mathrm b}$
  in~\eqref{eq:pbkl} is $|S_{\mathrm b}|$, the held-out bound subset, not the
  size of the underlying dataset.
\item \textbf{Joint confidence.} The Monte-Carlo and PAC-Bayes steps use separate
  failure probabilities $\delta_{\mathrm{MC}}=\delta_{\mathrm{PAC}}=0.005$,
  combined by a union bound to a reported confidence of $0.99$, rather than reusing
  one $\delta$ across both inversions.
\item \textbf{Monte-Carlo aggregation.} The empirical Gibbs risk is accumulated
  example-weighted over the bound set and averaged over exactly $m$ samples,
  fixing an off-by-one in the original per-batch averaging.
\end{enumerate}
The fourth affects reproducibility; the fifth affects whether the learnt-prior
construction satisfies the theorem's independence condition at all, though
neither changes the certificate arithmetic once the split is fixed:
\begin{enumerate}\setcounter{enumi}{3}
\item \textbf{Seed control.} Five independent seed streams (split, model, loader,
  Monte-Carlo, evaluation) are derived from one base seed, so the seed genuinely
  changes the split and the model. The evaluation stream is separate so that the
  descriptive test errors, which also sample posterior weights, do not inherit
  the RNG state the certificate's $m$ draws left behind: otherwise raising the
  Monte-Carlo budget would move a test error that does not depend on it.
\item \textbf{Label-independent partition.} The subset $S$ and its split into
  $S_0$ and $S_{\mathrm b}$ are drawn by a uniform random permutation rather than
  by taking a prefix. We first replaced the prefix split with a class-stratified
  one, then rejected that. Its class quotas are taken over the whole of $S$, which
  contains $S_{\mathrm b}$, so $S_0$ comes to depend on bound-set labels; and a
  sample with fixed class quotas is not the i.i.d.\ draw the bound assumes.
\end{enumerate}
All five rest on the explicit, saved bookkeeping of
Section~\ref{sec:compute}, which also keeps the six notions of error distinct:
the raw Monte-Carlo $0$--$1$ estimate, the corrected empirical Gibbs risk, the
certificate, and the stochastic, posterior-mean and ensemble test errors, so the
certified quantity is never silently conflated with a test error.

Each test targets a specific correction rather than basic pipeline execution, and
fixes a concrete input and an expected output rather than a plausible range: the
split tests assert the counts $6{,}000$ and $3{,}000$ at a tenth of the data, the
aggregation test asserts agreement to $10^{-9}$ between a single batch and a
ragged batching of a deterministic posterior, and the confidence test asserts
$0.99$ to $10^{-12}$ rather than $1-\delta$. Appendix~\ref{app:tests} lists all
forty-two with their tolerances.
A test that never fails defends nothing, so we checked that each bites rather
than assuming it: disabling the prior-selection union in the source makes two of
its five tests fail, and letting one $S_0$ example into the bound loader makes
two of the loader tests fail. They check the bookkeeping around the bound, not the bound itself.

\subsection{Training protocol}\label{sec:protocol}

Both the prior network and the posterior are trained by SGD with momentum and no
weight decay. The prior network minimises the negative log-likelihood of its
labels on $S_0$; the posterior minimises the objective of
Section~\ref{sec:objectives}. Table~\ref{tab:hparams} lists every setting, and one training configuration is
shared by every cell: its entries do
not vary with the dataset, the architecture, the objective or the prior family.
Only the evaluation batch cap differs, and for memory rather than training
reasons: the bound set is scored in a single batch for the FCN cells, in batches
of $2{,}500$ for the two CNN cells on $28\times28$ data, and $12{,}500$ on
CIFAR-10, so that the deeper networks' Monte-Carlo evaluation fits in device
memory. The cap changes how the estimate is batched, not what it estimates
(Section~\ref{sec:confidence}).

\begin{table}[h]\centering\small
\caption{The training configuration shared by every cell.}\label{tab:hparams}
\begin{tabular}{@{}llll@{}}
\toprule
Optimiser & SGD, no weight decay & Prior scale $\sigma_0$ & $0.03$ \\
Posterior learning rate & $5\times10^{-3}$ & Prior distribution & Gaussian \\
Posterior momentum & $0.95$ & $p_{\min}$ & $10^{-5}$ \\
Prior learning rate & $5\times10^{-3}$ & $\delta_{\mathrm{train}}$ & $0.025$ \\
Prior momentum & $0.99$ & $\delta_{\mathrm{MC}}$, $\delta_{\mathrm{PAC}}$ & $0.005$ each \\
Prior dropout & $0.2$ & Monte-Carlo samples $m$ & $2{,}000$ \\
Batch size & $250$ & Ensemble samples $L$ & $100$ \\
Prior epochs & $70$ & Prior/bound split & $0.5$ \\
Posterior epochs & $100$ & KL penalty $\beta$ (all objectives) & $0.1$ \\
\bottomrule
\end{tabular}
\end{table}

Since one fixed configuration replaces the reference's per-cell grid search, two
of these values were not simply copied.
The learning rates, momenta, $\sigma_0$, $p_{\min}$, batch size and
$\delta_{\mathrm{train}}$ are those the reference reports for MNIST, and the epoch
counts follow its statement that the networks converge by roughly $70$ epochs. One
value had to be recovered rather than copied: the reference implementation's
example script ships a posterior learning rate of $10^{-3}$, at which our first
grid gave a stochastic test error about four times the published one, whereas the
reference's text identifies $10^{-3}$ as converging slowly and $5\times10^{-3}$ as
the working value. We resolved that in favour of the text, but what prompted us
to look was a test-set discrepancy, so the learning-rate choice was indirectly
informed by the test set. It costs the certificate nothing, since it acts only on
the posterior, but it does mean the reported test errors are not a wholly
untouched estimate. Apart from that, the only sweep we ran ourselves was over the prior scale,
$\sigma_0\in\{0.01,0.02,0.05\}$ on MNIST at a reduced budget: the certificate was
almost unchanged across $[0.02,0.05]$ and degraded at $0.01$, so $\sigma_0=0.03$
sits in a flat region rather than at a tuned optimum. Those comparisons read
certificates computed on the bound set, so the prior was selected using
$S_{\mathrm b}$, and every certificate we report is corrected for that by the
union of Section~\ref{sec:confidence}. The reference grid-searches the same
parameter per cell without any such correction. This remains an explicitly
\emph{reduced-compute protocol}: the reference grid-searches learning rate, prior
scale and dropout per cell. Section~\ref{sec:mc} measures the effect of the
smaller $m$; the absent grid search is not quantified and is carried as a threat
to validity (Section~\ref{sec:evaluation}).

\subsection{Seeds and statistics}\label{sec:seeds}

Each headline cell is run over five base seeds ($0$--$4$), each genuinely
changing the data split, the model initialisation, the loader order and the
Monte-Carlo draws. Main-table quantities are means~$\pm$~standard deviations over
seeds, with per-seed values in the appendix. The across-seed variability of the
training procedure is not a PAC-Bayes confidence interval.

\subsection{Sensitivity and controlled-difficulty studies}\label{sec:studies}

Three further studies use the saved checkpoints. The \textbf{Monte-Carlo
sensitivity} curve isolates the effect of the sample budget $m$ on a fixed
posterior. At each budget
$m\in\{1{,}000,2{,}000,5{,}000,10{,}000,50{,}000,150{,}000\}$, up to the
reference's, we genuinely re-draw the Monte-Carlo estimate on the checkpoint and
recompute the certificate. Since the raw mean is an unbiased estimate of the
Gibbs risk, the only systematic route by which $m$ enters is the finite-MC slack
$\log(2/\delta_{\mathrm{MC}})/m$ of the first KL inversion, so we also compute an
analytical curve with the estimate held at its $m=1{,}000$ value; the two
agreeing is what shows the movement is slack and not a shifting estimate. The \textbf{small-data} sweep retrains the MNIST
FCN $f_{\text{quad}}$ learnt-prior cell on $100$, $50$, $20$ and $10\%$ of the
training data, with the corrected $n_{\mathrm b}$. The \textbf{controlled-difficulty} study
applies label noise $\eta\in\{0,0.1,0.2,0.4\}$ to the same cell with
architecture, sample size and compute held fixed.

The noise is symmetric: each example of $S$ is independently marked with
probability $\eta$ and a marked label is replaced by one drawn uniformly from the
nine \emph{other} classes, so a corrupted label is never the original and the
expected fraction of wrong labels is $\eta$. It is applied once, to $S$ as a whole
and before the prior/bound split, so $S_0$ and $S_{\mathrm b}$ are corrupted at
the same rate and the prior network also trains on noisy labels; the test set is
untouched, and the noise draw has its own seed. The certificate then bounds the
Gibbs risk under the
noisy-label distribution, whereas the test errors in the same rows are measured
against clean labels, so the two are not comparable. All three studies use the
same five seeds as the headline matrix.

\subsection{Compute and reproducibility}\label{sec:compute}

Experiments run on two NVIDIA RTX 3080 cards (20\,GB each) under PyTorch 2.5.1
and CUDA 12.4 \parencite{paszke2019}, and the full study comprises $115$ runs. Results are
aggregated from the per-run records into a single processed registry, which
de-duplicates by the full configuration and records the status and any exclusion
reason of every run, so that failed or non-finite runs are dropped with a stated
reason rather than silently. One criterion beyond the mechanical ones is the only
exclusion any run triggers. A
learnt-prior cell whose prior network never leaves the chance-level plateau is
not an instance of the protocol, since the posterior is then initialised at an
untrained prior, so we exclude a learnt-prior run whose prior network finishes
within five percentage points of chance, which for ten classes is $0.90$, on
$S_0$, the subset it was trained on. Scoring the screen there rather than on the
test set keeps it a training diagnostic: it consults no held-out data and exists
before any certificate does, so it cannot be tuned against a result. Separation
is wide: prior networks that trained reach at most $0.228$ on $S_0$, and the two
that did not sit at $0.866$ and $0.899$. Every table and figure reporting
experimental results in this dissertation is produced from that registry;
Figures~\ref{fig:chain}--\ref{fig:audit} are explanatory schematics rather than
results figures, and the pre-correction outputs are retained separately as an
archival audit record and are not used in any reported result.

This provenance is what allows each reported certificate to be recomputed and
checked independently. Each run stores all of its
seeds, the sample accounting ($n_{\mathrm{total}}$, $n_{\mathrm{selected}}$,
$n_{\mathrm{prior}}$, $n_{\mathrm{posterior}}$, $n_{\mathrm b}$), the three
confidence parameters, the raw Monte-Carlo estimate, the KL and the resulting
certificate, alongside the library versions, the hardware and the software version.
Any certificate can therefore be recomputed without re-training, which is how the
prior-selection union of Section~\ref{sec:confidence} was applied to all runs
after the fact. The registry carries the corrected and uncorrected values side by
side, and recomputing
the uncorrected column from the stored ingredients reproduces all $115$ recorded
certificates to $10^{-9}$.

The artefact is laid out so that each of those claims can be checked from a
command. One YAML file per headline cell fixes every entry of
Table~\ref{tab:hparams}, and the published values used in Table~\ref{tab:repro}
are held in a separate reference file annotated with the source column each was
read from, so the comparison cannot drift onto a different column of the
reference table. Appendix~\ref{app:repro} gives the directory layout, the
commands that regenerate a run, the matrix, the tests and every derived output,
and the pinned environment.
